\chapter*{Evidence and source boundaries}
The XOR code, framing trace, finite constrained-code implementation and independent
tests are teaching references, not a new storage architecture. Enumerative coding,
error correction and molecular storage have substantial prior art.
The raw byte archive is intentionally not sequence-constrained. The rank/unrank
code proves only its GC and run bounds. No primer validation, wet-lab recovery,
channel calibration or optimized production decoder is claimed.

Primary records were checked on 6 October 2026. Goldman and Organick were inspected
at publisher abstract/metadata depth (Organick's accessible primer-design captions
were also inspected). Chen et al. were read at abstract, introduction and selected
bias/model sections in the primary PMC article. Gimpel et al. were read at abstract
and selected introduction/initial error-analysis sections. Bar-Lev et al. were
inspected at abstract/data-code-availability depth; their correction explicitly
adds missing prior work to simulated-data methods. Li et al. were inspected at
abstract and selected codebook/overview/comparison sections, not all supplementary
proofs or experimental data. None of these studies was independently reproduced.

\begin{thebibliography}{99}
\bibitem{goldman} N. Goldman et al. Towards practical, high-capacity,
low-maintenance information storage in synthesized DNA.
\emph{Nature} 494, 77--80 (2013).
\url{https://doi.org/10.1038/nature11875}.
\bibitem{organick} L. Organick et al. Random access in large-scale DNA data storage.
\emph{Nature Biotechnology} 36, 242--248 (2018), corrected publisher record.
\url{https://doi.org/10.1038/nbt.4079}.
\bibitem{bias} Y.-J. Chen et al. Quantifying molecular bias in DNA data storage.
\emph{Nature Communications} 11, 3264 (2020).
\url{https://doi.org/10.1038/s41467-020-16958-3}.
\bibitem{twin} A. L. Gimpel, W. J. Stark, R. Heckel and R. N. Grass.
A digital twin for DNA data storage based on comprehensive quantification of errors
and biases. \emph{Nature Communications} 14, 6026 (2023).
\url{https://doi.org/10.1038/s41467-023-41729-1}.
\bibitem{deep-storage} D. Bar-Lev et al. Scalable and robust DNA-based storage via
coding theory and deep learning. \emph{Nature Machine Intelligence} 7, 639--649 (2025).
\url{https://doi.org/10.1038/s42256-025-01003-z}.
\bibitem{deep-correction} D. Bar-Lev et al. Author Correction to the preceding work.
Published 30 December 2025; \emph{Nature Machine Intelligence} 8, 134 (2026).
\url{https://doi.org/10.1038/s42256-025-01175-8}.
\bibitem{bijective} J. Li, M. Dong, Z. Qi and H. Guan.
Towards density-optimal DNA storage with practical bijective coding.
\emph{Nature Communications} 17, 10261 (2026).
\url{https://doi.org/10.1038/s41467-026-77024-y}.
\end{thebibliography}
